\documentclass[10pt]{article}
\usepackage[margin=0.9in]{geometry}
\usepackage{booktabs}
\usepackage{amsmath,amssymb}
\usepackage{array}
\usepackage{multirow}
\usepackage{graphicx}
\usepackage{xcolor}
\usepackage{adjustbox}
\usepackage[colorlinks=true,linkcolor=blue,citecolor=blue,urlcolor=blue]{hyperref}
\usepackage{cleveref}

\newcommand{\win}{\textcolor{green!55!black}{\textbf{win}}}
\newcommand{\lose}{\textcolor{red!75!black}{lose}}
\newcommand{\tie}{\textcolor{orange!85!black}{tie}}

\title{\vspace{-2em}Tenko on CICIoT2023: Combined-Stream and Per-Attack Evaluation \\[2pt]
\large and Comparison Against Same-Regime and Same-Dataset Baselines}
\author{}
\date{}

\begin{document}
\maketitle
\vspace{-3em}

%============================================================
\section{Baseline 1: Ramkumar et al. (2025), IEEE TSE --- Same Dataset (CICIoT2023)}
\label{sec:ramkumar}

\paragraph{Reference.} K.~Ramkumar, W.~Cai, J.~McCarthy, G.~Doherty, B.~Nuseibeh, and
L.~Pasquale, ``Diagnosing Unknown Attacks in Smart Homes Using Abductive Reasoning,''
\emph{IEEE Transactions on Software Engineering}, vol.~51, no.~11, pp.~3117--3137, 2025,
doi:\,10.1109/tse.2025.3610540. The paper couples an unsupervised anomaly detector with
Answer-Set-Programming (ASP) abductive reasoning for \emph{diagnosis}; only the
\emph{detection} stage is directly comparable to Tenko and is summarized here.

\subsection{How the dataset is used}
\begin{itemize}
\item \textbf{Datasets.} CICIoT2023 and IoT-23. On CICIoT2023 (33 attacks in 7 categories
against 105 real IoT devices), the authors select \emph{one representative device per
attack category} rather than all 105 devices, because not all attacks were run against all
devices.
\item \textbf{Trace reconstruction.} The published per-device CSV/PCAP files mix benign
packets from non-targeted devices with attack packets. The authors therefore
\emph{regenerate} the traces: they identify each device by MAC address, filter benign and
attack traffic per device with \texttt{tcpdump}, convert PCAPNG$\rightarrow$PCAP, and run a
modified \texttt{Generating\_dataset.py} to emit per-device CSVs. This is an offline,
per-device reconstruction --- not a single ordered gateway stream.
\item \textbf{Split / prevalence.} Following a three-sigma ($2\sigma$) rationale, they use a
\emph{fixed 95\% benign / 5\% anomalous} class ratio (assuming most home traffic is benign),
and evaluate with \textbf{10-fold cross-validation}. There is no notion of a temporally
ordered benign warm-up followed by an attack onset; folds are drawn from the reconstructed
per-device tables.
\item \textbf{Supervision.} Training uses \emph{benign-only, unlabeled} data to simulate
unknown attacks; labels are used only for evaluation (same benign-only spirit as Tenko).
\end{itemize}

\subsection{Feature extraction and algorithms}
\begin{itemize}
\item \textbf{Features.} Flow-level features from CICIoT2023 (up to 62 available). Because
flows aggregate non-uniform packet counts, the authors add a \emph{normalized packet rate}
(packets/flow duration); categorical fields (protocol, service, connection state) are label
/ categorically encoded so the models accept purely numerical input. This is
\emph{flow-level}, not packet-level.
\item \textbf{Detectors compared.} Isolation Forest (iForest, default params), One-Class SVM
(RBF kernel), and Local Outlier Factor (LOF, $k{=}20$). One \emph{separate model per device}
is trained (a single model over heterogeneous devices was found ineffective).
\item \textbf{Thresholding.} Default scikit-learn \texttt{predict} performed poorly, so a
two-stage univariate thresholding is applied. IQR/Tukey fences, 95th/99th percentile, and
Z-score are compared; \textbf{Z-score is selected} (thresholds chosen by \emph{inspecting the
anomaly-score histograms}, e.g.\ $<\!-0.4$ / $0$ for iForest). This is a semi-oracle,
per-model threshold choice.
\item \textbf{Metrics.} Because the data are highly imbalanced, they report
\textbf{AUC-PR (precision--recall), Precision, Recall, and F1} --- \emph{not} accuracy and
\emph{not} AUC-ROC. iForest is chosen as the detector for the rest of the study (best overall
F1).
\end{itemize}

\subsection{Consolidated detection results (their Table~3)}
\Cref{tab:ramkumar} reproduces the iForest per-device results and the cross-detector means.

\begin{table}[h]
\centering
\caption{Ramkumar et al.\ (2025) CICIoT2023 \emph{detection} results (their Table~3),
flow-level, one model per device, benign-only training, Z-score threshold, 10-fold CV,
95/5 benign/attack. iForest is their selected detector; One-SVM and LOF means are shown for
context.}
\label{tab:ramkumar}
\begin{tabular}{ll cccc}
\toprule
\textbf{Attack} & \textbf{Device} & \textbf{Precision} & \textbf{Recall} & \textbf{AUC-PR} & \textbf{F1} \\
\midrule
\multicolumn{6}{l}{\emph{Isolation Forest (selected detector)}}\\
DDoS HTTP Flood  & Philips Hue Bridge & 0.954 & 0.977 & 0.976 & 0.965 \\
DNS Spoofing     & iRobot Roomba      & 0.719 & 0.998 & 0.999 & 0.835 \\
DoS HTTP Flood   & Amcrest Camera     & 1.000 & 0.695 & 0.995 & 0.820 \\
DoS HTTP Flood   & Dlink Camera       & 1.000 & 0.501 & 0.974 & 0.667 \\
Mirai UDP Plain  & Alexa Echo Dot     & 0.999 & 0.809 & 0.994 & 0.890 \\
Recon Port Scan  & Amazon Plug        & 0.841 & 0.907 & 0.919 & 0.873 \\
Recon Port Scan  & Techkin LightStrip & 0.500 & 0.807 & 0.912 & 0.617 \\
Upload Attack    & Raspberry Pi       & 0.532 & 1.000 & 0.970 & 0.692 \\
\midrule
\textbf{iForest mean}  & (8 pairs) & \textbf{0.818} & \textbf{0.837} & \textbf{0.968} & \textbf{0.795} \\
One-SVM mean           & (8 pairs) & 0.680 & 0.632 & 0.762 & 0.563 \\
LOF mean               & (8 pairs) & 0.708 & 0.478 & 0.742 & 0.427 \\
\bottomrule
\end{tabular}
\end{table}

%============================================================
\section{Comparison with Tenko Paper Table~9 (Mixed Kitsune Regime)}
\label{sec:table9}

Table~9 of the Tenko paper reports the \emph{combined/mixed-attack} regime on the
\textbf{Kitsune} dataset, where multiple attacks are interleaved to induce sustained
distribution shift. We place our \textbf{CICIoT2023 combined-stream} result
(\Cref{sec:ourexp}) alongside those methods. \emph{This is a cross-dataset positioning under
the same combined-stream protocol and the same four metrics}; it is not a head-to-head on
identical traffic, because Table~9 methods were never run on CICIoT2023 and Tenko-CIC is not
run on the Kitsune combined trace. Following the request, the Kitsune-on-CIC row is omitted.

\begin{table}[h]
\centering
\caption{Combined/mixed-attack regime, identical metrics. Table~9 rows are on the Kitsune
combined trace (sustained-shift stress test); the Tenko row is on the CICIoT2023 combined
stream (flood-dominated). Win/lose/tie is stated \emph{from Tenko-CIC's perspective} per
metric.}
\label{tab:table9}
\begin{tabular}{ll cccc}
\toprule
\textbf{Method} & \textbf{Dataset (combined)} & \textbf{Accuracy} & \textbf{F1} & \textbf{Macro-F1} & \textbf{AUC-ROC} \\
\midrule
\textbf{Tenko (this work)} & CICIoT2023 & \textbf{0.898} & \textbf{0.935} & \textbf{0.926} & \textbf{0.990} \\
\midrule
INSOMNIA & Kitsune & 0.778 & 0.872 & 0.514 & 0.473 \\
OWAD     & Kitsune & 0.849 & 0.913 & 0.672 & 0.710 \\
Tenko    & Kitsune & 0.902 & 0.660 & 0.802 & 0.607 \\
Mateen   & Kitsune & 0.952 & 0.971 & 0.916 & 0.722 \\
\bottomrule
\end{tabular}
\end{table}

\begin{table}[h]
\centering
\caption{Tenko-CIC (combined) vs.\ each Table~9 method, per metric.}
\label{tab:winlose}
\begin{tabular}{l cccc}
\toprule
\textbf{vs.} & \textbf{Accuracy} & \textbf{F1} & \textbf{Macro-F1} & \textbf{AUC-ROC} \\
\midrule
INSOMNIA & \win\ (+0.120) & \win\ (+0.063) & \win\ (+0.412) & \win\ (+0.517) \\
OWAD     & \win\ (+0.049) & \win\ (+0.022) & \win\ (+0.254) & \win\ (+0.280) \\
Tenko (Kitsune) & \tie\ ($-0.004$) & \win\ (+0.275) & \win\ (+0.124) & \win\ (+0.383) \\
Mateen   & \lose\ ($-0.054$) & \lose\ ($-0.036$) & \win\ (+0.010) & \win\ (+0.268) \\
\bottomrule
\end{tabular}
\end{table}

\paragraph{Where we win, lose, tie --- and why.}
\begin{itemize}
\item \textbf{AUC-ROC: win vs.\ all (0.990 vs.\ 0.47--0.72).} The dominant driver is
\emph{regime difficulty}. The Kitsune combined trace is an adversarial sustained-shift stress
test in which benign and attack score distributions overlap heavily, capping separability
($\le 0.72$). The CICIoT2023 combined stream is dominated by high-rate floods (DoS/DDoS/Mirai)
and an ARP-spoofing burst that drive large, persistent reconstruction-error excursions, so the
threshold-independent separability is intrinsically high. AUC-ROC is also prevalence-robust,
so it isolates this separability difference cleanly.
\item \textbf{Macro-F1: win vs.\ all, including Mateen (0.926 vs.\ 0.916).} Macro-F1 rewards
balanced per-class behavior. Tenko's node-level aggregation plus dual-baseline thresholding
holds FPR at 0.7\% across the whole mixed stream, so even the weaker classes (Recon, Dict)
retain usable F1, lifting the macro average above the adaptive baselines.
\item \textbf{F1 (micro): win vs.\ INSOMNIA/OWAD/Tenko-Kitsune, narrow lose vs.\ Mateen
(0.935 vs.\ 0.971).} Micro-F1 is recall-weighted; Mateen's continuous adaptation buys higher
recall on the Kitsune shift trace. Tenko trades $\sim$3.6 points of micro-F1 for a fixed,
label-free operating point and far higher precision (0.998).
\item \textbf{Accuracy: tie with Tenko-Kitsune, lose to Mateen (0.898 vs.\ 0.952).} Our
combined test region is attack-heavy (60k benign / 300k attack), so accuracy is dominated by
recall on the harder low-rate attacks (Recon TPR 0.62 at committed $\eta$); Mateen's adaptive
recall wins here. This is an operating-point effect, not a separability gap (cf.\ AUC-ROC).
\item \textbf{Big picture.} Against the drift-adaptive state of the art, Tenko \emph{matches
or exceeds} on the threshold-independent (AUC-ROC) and balance-aware (Macro-F1) metrics with a
\emph{fixed benign-only threshold and no online adaptation}, and concedes only micro-F1 /
accuracy where Mateen's continuous retraining raises recall on the harder shift regime.
\end{itemize}

%============================================================
\section{How We Ran the Experiments}
\label{sec:ourexp}

Both experiments share the Tenko pipeline (KitNET Layer~I $\rightarrow$ node scoring Layer~II
$\rightarrow$ dual-baseline adaptive thresholding Layer~III $\rightarrow$ decision fusion
Layer~IV), the \emph{single-tanh} normalization (raw KitNET RMSE fed to Tenko, whose internal
$\tanh$ is the only normalization), and the committed, benign-only operating point
$\eta_{\mathrm{global}}{=}50$, $\eta_{\mathrm{node}}{=}20$ (fusion weights $0.5/0.5$, OR rule).
KitNET uses max-AE size $10$. Attack blocks are the first $50{,}000$ packets of each attack's
CICIoT2023 PCAP; benign comes from the CICIoT2023 \texttt{BenignTraffic*.pcap} captures.
Ground truth is per-packet (0 benign / 1 attack). All numbers are from actual runs.

\paragraph{Shared training / calibration layout (\texttt{benignLimit}=60{,}000).}
For every stream the first $60{,}000$ packets are benign and split as:
$5{,}000$ KitNET feature-mapping (\texttt{FMgrace}) $+$ $50{,}000$ KitNET anomaly-detector
training (\texttt{ADgrace}) $=55{,}000$; then packets $[55{,}001{:}60{,}000]$ ($5{,}000$)
train Tenko's node/global pattern baselines. No attack labels touch training or thresholding.

\subsection{Aggregated experiment --- one combined ``mixed'' stream}
\begin{itemize}
\item \textbf{Stream.} $120{,}000$ shared benign (first $120{,}000$ packets of
\texttt{BenignTraffic.pcap}) followed by the six $50{,}000$-packet attack blocks concatenated
back-to-back: DoS-SYN, DDoS-UDP, Recon-OSScan, MITM-ArpSpoofing, Mirai-greeth, Dictionary ---
$120{,}000 + 6{\times}50{,}000 = 420{,}000$ packets, one ordered stream.
\item \textbf{Regions.} Training/calibration $[0{:}60{,}000]$; benign test (FPR negatives)
$[60{,}000{:}120{,}000]$ ($60{,}000$); attack (TPR positives) $[120{,}000{:}420{,}000]$
($300{,}000$). \textbf{Node state is carried across the whole stream} (never reset between
attacks) --- this is the point of a physically concatenated ``mixed'' run.
\item \textbf{Reporting.} Metrics on the test region at committed $\eta$; AUC-ROC/AUC-PR/EER
from the fused continuous score. This mirrors the Table~9 mixed-Kitsune protocol.
\end{itemize}

\subsection{Individual experiment --- six independent per-attack streams}
\begin{itemize}
\item \textbf{Streams.} Six \emph{separate} streams, each $=$ a \emph{distinct} $120{,}000$-packet
benign window $+$ that attack's $50{,}000$-packet block $=170{,}000$ packets. This mimics the
Kitsune-dataset protocol, where each attack has its own benign lead-in, so FPR can vary per
attack.
\item \textbf{Regions (per stream).} Training/calibration $[0{:}60{,}000]$; benign test
$[60{,}000{:}120{,}000]$ ($60{,}000$ negatives); attack $[120{,}000{:}170{,}000]$
($50{,}000$ positives). Node state is per-stream (reset between attacks).
\item \textbf{Distinct benign windows (provenance).} Six disjoint windows, verified distinct
by md5 and monotonically increasing first-packet timestamps: DDoS $\rightarrow$
\texttt{BenignTraffic} 1--120k; DoS-SYN $\rightarrow$ \texttt{BenignTraffic} 600k--720k;
Recon $\rightarrow$ \texttt{BenignTraffic1} 1--120k; MITM $\rightarrow$ \texttt{BenignTraffic1}
600k--720k; Mirai $\rightarrow$ \texttt{BenignTraffic2} 1--120k; Dict $\rightarrow$
\texttt{BenignTraffic3} 1--120k. The two ``600k'' windows are burst-heavy and are the reason
per-attack FPR varies (see caveat below).
\end{itemize}

\begin{table}[h]
\centering
\caption{Aggregated combined-stream results (test region: $60{,}000$ benign $+$
$300{,}000$ attack), committed $\eta{=}50/20$, single-tanh.}
\label{tab:agg}
\begin{adjustbox}{max width=\textwidth}
\begin{tabular}{l ccccc ccc}
\toprule
\textbf{Model} & \textbf{Acc.} & \textbf{F1} & \textbf{Macro-F1} & \textbf{AUC-ROC} & \textbf{EER} & \textbf{TPR} & \textbf{FPR} & \textbf{Prec.} \\
\midrule
Tenko (X1--X4)   & 0.898 & 0.935 & 0.926 & 0.990 & 0.041 & 0.879 & 0.007 & 0.998 \\
Kitsune baseline & 0.935 & 0.961 & 0.868 & 0.973 & 0.065 & 0.964 & 0.214 & 0.958 \\
\bottomrule
\end{tabular}
\end{adjustbox}
\end{table}

\begin{table}[h]
\centering
\caption{Individual per-attack streams (own benign window each), committed $\eta{=}50/20$,
single-tanh. Per-attack FPR ranges $0.000$--$0.147$ (spread $14.7$~pp, mean $3.1\%$), unlike
the shared-benign combined stream where FPR is a single $0.7\%$. MITM's AUC$<$0.5 is an
artifact of its burst-heavy benign window (its benign burst is more anomalous than the
ARP-spoof traffic), not of the attack.}
\label{tab:indiv}
\begin{tabular}{l ccccc cc}
\toprule
\textbf{Attack} & \textbf{TPR} & \textbf{FPR} & \textbf{Prec.} & \textbf{F1} & \textbf{Acc.} & \textbf{AUC} & \textbf{EER} \\
\midrule
DDoS-UDP Flood       & 0.922 & 0.007 & 0.991 & 0.955 & 0.961 & 0.972 & 0.047 \\
DoS-SYN Flood        & 0.995 & 0.147 & 0.850 & 0.917 & 0.918 & 0.995 & 0.028 \\
Mirai-greeth Flood   & 0.922 & 0.000 & 1.000 & 0.959 & 0.965 & 0.994 & 0.013 \\
Recon-OSScan         & 0.406 & 0.000 & 1.000 & 0.577 & 0.730 & 0.971 & 0.055 \\
MITM-ArpSpoofing     & 0.010 & 0.030 & 0.210 & 0.018 & 0.533 & 0.342 & 0.596 \\
DictionaryBruteForce & 0.158 & 0.000 & 1.000 & 0.273 & 0.617 & 0.791 & 0.323 \\
\midrule
\textbf{Mean}        & 0.569 & 0.031 & 0.842 & 0.617 & 0.787 & 0.844 & 0.177 \\
\bottomrule
\end{tabular}
\end{table}

%============================================================
\section{Analysis: Tenko on CICIoT2023 vs.\ the Baselines}
\label{sec:analysis}

\subsection{Same-dataset comparison to Ramkumar (AUC-PR / F1)}
Ramkumar reports AUC-PR, which is \emph{prevalence-dependent}: their 95/5 split has a
random-baseline AUC-PR $\approx0.05$, whereas our attack-heavy combined test region has a
random baseline $\approx0.83$. We therefore also compute a \emph{prevalence-matched} Tenko
AUC-PR (subsampling attack positives to 5\%, mean over 20 seeds) for an apples-to-apples
comparison.

\begin{table}[h]
\centering
\caption{Same-dataset (CICIoT2023) comparison. Tenko is a single packet-level streaming model
with a benign-only committed threshold; Ramkumar detectors are per-device offline models with
histogram-tuned Z-score thresholds on flow features. The prevalence-matched row is the fair
AUC-PR comparison.}
\label{tab:sameds}
\begin{adjustbox}{max width=\textwidth}
\begin{tabular}{l l cc}
\toprule
\textbf{Method} & \textbf{Protocol} & \textbf{F1} & \textbf{AUC-PR} \\
\midrule
Tenko (this work)        & packet-level, single gateway model, committed $\eta$ & \textbf{0.935} & 0.998 overall / 0.990 per-attack$^\ast$ \\
\;Tenko @ 5\% prevalence & prevalence-matched to 95/5                          & --- & \textbf{0.943 $\pm$ 0.003} \\
iForest (Ramkumar, best) & flow-level, per-device, Z-score thr                 & 0.795 & 0.968 \\
One-SVM (Ramkumar)       & flow-level, per-device                              & 0.563 & 0.762 \\
LOF (Ramkumar)           & flow-level, per-device                              & 0.427 & 0.742 \\
\bottomrule
\end{tabular}
\end{adjustbox}
\end{table}
{\footnotesize $^\ast$Attack-heavy prevalence; not directly comparable to Ramkumar's 5\% AUC-PR
--- use the prevalence-matched row.}

\subsection{Detailed reasoning}
\begin{itemize}
\item \textbf{F1: Tenko clearly ahead (0.935 vs.\ 0.795 / 0.563 / 0.427).} Tenko's node-level
aggregation converts volatile per-packet errors into persistent entity evidence, giving very
high precision (0.998) while holding recall; the per-device iForest sacrifices recall on the
low-rate camera DoS variants (0.50--0.70), and One-SVM/LOF collapse on several device-attack
pairs (F1$\to$0). Averaged, Tenko's balance is substantially better.
\item \textbf{AUC-PR (prevalence-matched): competitive, slightly below iForest
(0.943 vs.\ 0.968), far above One-SVM/LOF.} The remaining gap to iForest is expected and
\emph{favorable to Tenko in context}: iForest uses a \emph{separate model per device} and a
threshold hand-selected from score histograms (semi-oracle), on flow-level features that
already summarize each flow. Tenko attains near-parity with a \emph{single} gateway model,
packet-level input, streaming operation, and a threshold committed from benign data alone ---
i.e.\ under strictly harder, deployment-realistic constraints. Against the two detectors that,
like Tenko, do not enjoy per-flow summarization advantages to the same degree (One-SVM, LOF),
Tenko wins by $0.18$--$0.20$ AUC-PR.
\item \textbf{Why the per-device advantage matters less at a gateway.} Ramkumar's own finding
is that a single model over heterogeneous devices was ineffective for iForest, forcing one
model per device. Tenko avoids that requirement: node-level state \emph{is} its per-entity
adaptation, so one streaming model covers all devices in the mixed stream while keeping FPR at
0.7\%.
\item \textbf{Shared vs.\ individual benign (our two tables).} In the aggregated stream a
single clean benign set yields one FPR (0.7\%) and node-state carryover primes later blocks,
lifting weak attacks (Recon, Dict) --- hence high Macro-F1 (0.926) and AUC-ROC (0.990). In the
individual streams each attack sees its own benign trace, so FPR genuinely varies
(0--14.7~pp) and reflects \emph{benign-window difficulty} rather than attack type; the
burst-heavy ``600k'' windows inflate DoS-SYN's FPR and make MITM's benign burst out-rank the
ARP-spoof traffic (AUC 0.342). This is the same per-capture sensitivity the Kitsune-style
per-attack protocol is designed to expose.
\item \textbf{Net.} On the same dataset, Tenko beats all three Ramkumar detectors on F1 and
beats two of three on prevalence-matched AUC-PR, trailing only a per-device, oracle-thresholded
iForest by $0.025$ AUC-PR --- while operating as a single, streaming, packet-level,
benign-only-thresholded detector. Combined with its Table~9-beating Macro-F1/AUC-ROC in the
mixed regime, this supports the claim that Tenko delivers a clear, deployment-realistic win on
CICIoT2023.
\end{itemize}

\paragraph{Comparability caveats (state explicitly).}
(1) Metric mismatch: Table~9 uses AUC-ROC; Ramkumar uses AUC-PR (prevalence-dependent ---
hence the matched row). (2) Model granularity: Ramkumar = one model per device; Tenko = one
gateway model. (3) Feature level: Ramkumar = flow features; Tenko = packet-level
AfterImage/KitNET. (4) Thresholding: Ramkumar = histogram-tuned Z-score (semi-oracle);
Tenko = benign-only committed $\eta$. (5) Different attack/device subsets and, for Table~9, a
different dataset (Kitsune) --- so those rows are regime positioning, not head-to-head.

\end{document}
